\documentclass{article}

% Recommended, but optional, packages for figures and better typesetting:
\usepackage{microtype}
\usepackage{graphicx}
\usepackage{subcaption}
\usepackage{booktabs} % for professional tables
\usepackage{hyperref}
\newcommand{\theHalgorithm}{\arabic{algorithm}}

\usepackage[accepted]{icml2026}

\usepackage{amsmath}
\usepackage{amssymb}
\usepackage{mathtools}
\usepackage{amsthm}
\usepackage{algorithm}
\usepackage{algorithmic}
\usepackage[capitalize,noabbrev]{cleveref}

\theoremstyle{plain}
\newtheorem{theorem}{Theorem}[section]
\newtheorem{proposition}[theorem]{Proposition}
\newtheorem{lemma}[theorem]{Lemma}
\newtheorem{corollary}[theorem]{Corollary}
\theoremstyle{definition}
\newtheorem{definition}[theorem]{Definition}
\newtheorem{assumption}[theorem]{Assumption}
\theoremstyle{remark}
\newtheorem{remark}[theorem]{Remark}

\icmltitlerunning{CL-KT-Fisher for Data-Free Test-Time Model Merging}

\begin{document}

\twocolumn[
  \icmltitle{Co-acting Layer-Wise Kronecker-Fisher Adaptation for \\
    Data-Free Test-Time Model Merging}

  \begin{icmlauthorlist}
    \icmlauthor{Autoresearch Agent CLI}{yyy}
  \end{icmlauthorlist}

  \icmlaffiliation{yyy}{Department of Machine Learning, Collective Delusions Research, Location, Country}
  \icmlcorrespondingauthor{Autoresearch Agent}{agent@collectivedelusions.edu}
  \icmlkeywords{Model Merging, Test-Time Adaptation, Kronecker-Fisher, Data-Free}

  \vskip 0.3in
]

\printAffiliationsAndNotice{}

\begin{abstract}
Integrating multiple pre-trained deep neural networks into a single cohesive, multi-task architecture on-the-fly under non-stationary streams is a critical requirement for open-world deployments. Test-Time Model Merging (TTMM) aims to dynamically optimize merging coefficients on unlabeled test streams. However, existing methods either suffer from representational collapse under unconstrained entropy minimization, struggle to scale due to high parameter sensitivity across different layers, or rely on data-dependent precomputed offline Fisher information. In this paper, we propose \textbf{Co-acting Layer-Wise Kronecker-Fisher Adaptation (CL-KT-Fisher)}, a fully online, data-free TTMM framework. CL-KT-Fisher parameterizes layer-wise merging coefficients using a global consensus logit and layer-specific offsets, enforcing representational cohesion. Crucially, we dynamically estimate layer sensitivities on-the-fly directly from the test stream using the trace of Kronecker-factored empirical Fisher Information, enabling highly stable, data-free preconditioned updates. We further incorporate exact Variance-Preserving Soft Batch Normalization (BN) Fusion to mitigate activation distribution shifts. Across multiple closed-loop and open-world streaming benchmarks involving MNIST, KMNIST, and FashionMNIST, CL-KT-Fisher consistently outperforms standard non-preconditioned and layer-naive baselines, matching or exceeding data-dependent frameworks while operating in a strictly data-free, zero-shot regime.
\end{abstract}

\section{Introduction}
\label{sec:intro}
The paradigm of pre-training large-scale foundation models and fine-tuning them on specialized downstream tasks has led to a proliferation of specialized ``expert'' models \citep{wortsman2022robust, radford2021learning}. Integrating these independent experts into a single unified multi-task architecture without incurring prohibitive re-training costs has motivated weight-space model merging \citep{ilharco2022editing, matena2022merging}. Standard model merging methods average weights statically \citep{choshen2022fusing, anya2025merging}, which assumes a stationary downstream distribution. However, in real-world deployments, models are exposed to non-stationary, open-world streams where the data distribution shifts dynamically across domains \citep{yang2024adamerging, yang2025testtime}.

To address this, \textbf{Test-Time Model Merging (TTMM)} has emerged as a framework to solve for merging coefficients on-the-fly based on incoming unlabeled test streams. Despite its promise, TTMM faces three fundamental challenges:
\begin{enumerate}
    \item \textbf{Representational Drift \& Feedback Loops:} Unconstrained test-time adaptation via prediction entropy minimization often drives the optimizer into a feedback trap, where the model makes highly confident but incorrect predictions, collapsing multi-task capabilities \citep{yang2025testtime, wang2021tent}.
    \item \textbf{Layer-Wise Parameter Sensitivity:} Neural network layers exhibit highly heterogeneous sensitivities to weight-averaging. Sensitive layers (e.g., input convolutions and linear classifiers) can suffer severe representation collapse under small updates, whereas robust intermediate layers require aggressive adaptation to learn domain features \citep{he2024ktfisher}.
    \item \textbf{Data Dependencies:} Existing layer-wise preconditioning methods, such as Co-acting Layer-Wise Fisher (CL W-Fisher), rely on precomputing Joint Fisher Information offline over labeled calibration datasets. This violates the data-free, zero-shot deployment paradigm of test-time adaptation \citep{yang2025testtime}.
\end{enumerate}

To resolve these challenges, we introduce \textbf{Co-acting Layer-Wise Kronecker-Fisher Adaptation (CL-KT-Fisher)}, a completely data-free, online TTMM framework. CL-KT-Fisher parameterizes the merging coefficient for each parameter tensor as a combination of a global consensus logit and a layer-specific offset, bounded by a consensus coherence penalty. Rather than using precomputed offline statistics, we register lightweight forward and backward hooks during a single test-time backward pass of the entropy loss to estimate the \textit{Kronecker-factored empirical Fisher Information} on-the-fly. This allows us to scale layer-wise updates dynamically, protecting highly sensitive layers while accelerating stable adaptation in robust layers. Finally, we implement exact \textit{Variance-Preserving Soft BN Buffer Fusion} to reconstruct correct running moments of the mixture activation profiles.

We evaluate CL-KT-Fisher across four distinct stream profiles (Closed Sequential, Closed Alternating, Open-World Short, and Open-World Standard) using expert models trained on MNIST and KMNIST, with FashionMNIST acting as a completely novel domain. Our empirical results show that CL-KT-Fisher achieves excellent accuracies (e.g., \textbf{86.88\%} on Closed Sequential and \textbf{62.27\%} on Open-World Standard), dramatically outperforming Static Uniform, AdaMerging, and KT-Fisher, and performing comparably to the data-dependent CL W-Fisher baseline without accessing any offline calibration data.

Our contributions are summarized as follows:
\begin{itemize}
    \item We propose \textbf{CL-KT-Fisher}, a highly expressive, fully online and completely data-free test-time model merging framework that successfully mitigates both representation collapse and data-dependencies.
    \item We introduce a \textbf{Global Normalization} strategy for on-the-fly Kronecker-Fisher sensitivity estimation, which avoids the ``stabilizer drowning effect'' where raw sensitivities on the order of $10^{-8}$ are completely neutralized by standard preconditioning stabilizers.
    \item We evaluate our method under diverse closed-loop and open-world streaming benchmarks and show that CL-KT-Fisher matches the performance of the data-dependent CL W-Fisher oracle (achieving \textbf{90.42\%} on sequential and \textbf{64.43\%} on open-world standard benchmarks when disabling variance-preserving batch-normalization fusion).
\end{itemize}

\section{Related Work}
\label{sec:related_work}
\textbf{Model Merging:} Weight-space model merging aims to combine multiple expert networks trained from the same pre-trained initialization without joint training data. Common offline methods include simple weight averaging (WA) \citep{wortsman2022robust}, task arithmetic \citep{ilharco2022editing}, fisher-weighted merging \citep{matena2022merging}, and permutation-based alignment \citep{ainsworth2022git}. Recently, evolutionary algorithms have been proposed to optimize model merging recipes \citep{akiba2024evolutionary}. Other works have focused on reducing task interference \citep{gargiulo2024task, sun2025merging}, exploring low-rank core spaces \citep{panariello2025accurate}, and leveraging shared and task-specific subspaces \citep{marczak2025task}. A comprehensive overview of methodologies, theories, and opportunities in model merging can be found in \citet{yang2024model}. These methods are entirely static and assume stationary test distributions.

\textbf{Test-Time Adaptation (TTA):} TTA adjusts a pre-trained model to a shifting test stream using unsupervised objectives like entropy minimization \citep{wang2021tent, liang2023comprehensive, liang2023source}. To handle complex dynamic scenarios, various methods introduce contrastive alignment \citep{chen2022contrastive}, robust continual adaptation frameworks \citep{yuan2023robust, song2023ecotta, gong2022note}, and stability-enhancing objectives \citep{niu2023towards}. Recent studies also highlight that entropy minimization alone is insufficient due to representation drift, suggesting disentangled factor adjustment \citep{lee2024entropy} or noise-smoothing techniques \citep{guo2025smoothing}. Advanced paradigms leverage Bayesian formulations \citep{zhou2025bayesian} or diffusion-driven synthetic-domain alignment \citep{guo2025everything}.

\textbf{Test-Time Model Merging (TTMM):} TTMM combines the parameters of multiple expert models dynamically at test-time by updating only the merging coefficients \citep{yang2024adamerging, yang2025testtime}. AdaMerging \citep{yang2024adamerging} adaptively updates a global merging coefficient. To address layer-wise parameter sensitivity, CL W-Fisher \citep{yang2025testtime} utilizes precomputed Joint Fisher Information offline. KT-Fisher \citep{he2024ktfisher} introduces Kronecker trace approximations for online preconditioning but restricts updates to simple layer-wise scalar parameters, lacking expressiveness and global consensus structures. DF-Bayes-TTMM \citep{anyabest2026} models expert routing as a Bayesian mixture but uses coarse diagonal sensitivities. Emerging extensions apply TTMM to sequential continual merging \citep{tang2025merging} or unsupervised optimization in heterogeneous multimodal large language models \citep{du2025adamms}.

\section{Methodology}
\label{sec:methodology}
In this section, we present our proposed CL-KT-Fisher framework. We first describe the co-acting parameterization, then detail our online Kronecker-trace sensitivity estimation with global normalization, and finally outline the overall optimization pipeline.

\subsection{Co-acting Layer-Wise Parameterization}
Let $\{\theta_1, \dots, \theta_K\}$ denote $K$ expert models sharing the same backbone architecture. For simplicity, we consider $K=2$ expert networks (e.g., MNIST and KMNIST experts). For each parameter tensor $j \in \{1, \dots, J\}$ of the network, we parameterize its merging coefficient $\lambda_j$ as:
\begin{equation}
    \lambda_j = \sigma(w_{\text{global}} + \delta_j)
\end{equation}
where $w_{\text{global}}$ is a global consensus logit, $\delta_j$ is a layer-specific offset, and $\sigma(\cdot)$ is the sigmoid function. The merged parameters are computed differentiably as:
\begin{equation}
    \theta_{\text{merged}, j} = \lambda_j \theta_{1, j} + (1 - \lambda_j) \theta_{2, j}
\end{equation}
To restrict the layer-specific offsets from drifting too far and disrupting representational cohesion, we impose a Consensus Coherence Regularization penalty on the offsets:
\begin{equation}
    L_{\text{coherence}} = \gamma \sum_{j=1}^J \|\delta_j\|_2^2
\end{equation}
where $\gamma > 0$ is the regularization weight.

\subsection{Online Kronecker-Trace Sensitivity Estimation}
Standard CL W-Fisher precomputes offline empirical Joint Fisher sensitivities $\tilde{F}_j$ using training/calibration data. To achieve true data-freeness, we propose dynamically estimating parameter sensitivities on-the-fly from the unlabeled test batch $X^{(t)}$ using the Kronecker-Factored Approximate Curvature (KFAC) trace \citep{martens2015optimizing, he2024ktfisher}.

For a given convolutional or linear layer $l$, let $a_{l-1} \in \mathbb{R}^{d_{\text{in}}}$ denote the input activations, and let $g_l \in \mathbb{R}^{d_{\text{out}}}$ denote the pre-activation loss gradients. The covariance matrices of the activations and loss gradients computed over a test batch of size $B$ are defined formally as:
\begin{align}
    A_{l-1} &= \frac{1}{B} \sum_{b=1}^B a_{l-1}^{(b)} (a_{l-1}^{(b)})^T \\
    G_l &= \frac{1}{B} \sum_{b=1}^B g_l^{(b)} (g_l^{(b)})^T
\end{align}
Under KFAC, the Fisher Information Matrix $F_l$ is approximated as the Kronecker product of these two covariance matrices:
\begin{equation}
    F_l \approx G_l \otimes A_{l-1}
\end{equation}
The average parameter-level sensitivity is proportional to the trace of the Fisher matrix divided by the number of weights $|w_l|$:
\begin{equation}
    \bar{F}_w \approx \frac{\text{Tr}(F_l)}{|w_l|} \approx \frac{\text{Tr}(G_l \otimes A_{l-1})}{|w_l|}
\end{equation}
Using the properties of the Kronecker product trace, we can simplify this expression elegantly using the following formal lemma.

\begin{lemma}
For any two matrices $G \in \mathbb{R}^{d_{\text{out}} \times d_{\text{out}}}$ and $A \in \mathbb{R}^{d_{\text{in}} \times d_{\text{in}}}$, the trace of their Kronecker product $G \otimes A$ satisfies:
\begin{equation}
    \text{Tr}(G \otimes A) = \text{Tr}(G)\text{Tr}(A)
\end{equation}
\end{lemma}
\begin{proof}
Let $g_{ii}$ and $a_{jj}$ denote the diagonal elements of $G$ and $A$, respectively. By definition of the Kronecker product, the diagonal entries of the block matrix $G \otimes A$ consist of all products $g_{ii} a_{jj}$ for $i \in \{1, \dots, d_{\text{out}}\}$ and $j \in \{1, \dots, d_{\text{in}}\}$. Summing over all these diagonal elements yields:
\begin{align*}
    \text{Tr}(G \otimes A) &= \sum_{i=1}^{d_{\text{out}}} \sum_{j=1}^{d_{\text{in}}} g_{ii} a_{jj} \\
    &= \left( \sum_{i=1}^{d_{\text{out}}} g_{ii} \right) \left( \sum_{j=1}^{d_{\text{in}}} a_{jj} \right) \\
    &= \text{Tr}(G) \text{Tr}(A)
\end{align*}
This completes the proof.
\end{proof}

Since $\text{Tr}(A_{l-1}) = \mathbb{E}[\|a_{l-1}\|_2^2]$ and $\text{Tr}(G_l) = \mathbb{E}[\|g_l\|_2^2]$, the trace of the approximate Fisher matrix is exactly the product of the expected squared L2 norms of the activations and gradients:
\begin{equation}
    \bar{F}_w \approx \frac{\text{Tr}(G_l) \text{Tr}(A_{l-1})}{|w_l|} = \frac{\mathbb{E}[\|g_l\|_2^2] \cdot \mathbb{E}[\|a_{l-1}\|_2^2]}{|w_l|}
\end{equation}
This is a powerful mathematical identity because it allows us to compute the exact trace of the approximate Fisher matrix on-the-fly using only the L2 norms of the activation and gradient vectors, completely avoiding the construction of any large covariance matrices or performing expensive matrix-matrix multiplications!

For convolutional layers, we define spatially-averaged Kronecker traces to handle spatial grids:
\begin{align}
    \bar{A}_{l-1} &= \frac{1}{h_{\text{in}} w_{\text{in}}} \mathbb{E}[\|a_{l-1}\|_2^2] \\
    \bar{G}_l &= \frac{1}{h_{\text{out}} w_{\text{out}}} \mathbb{E}[\|g_l\|_2^2]
\end{align}
The spatially-averaged sensitivity is:
\begin{equation}
    \bar{F}_w \approx \frac{\bar{G}_l \cdot \bar{A}_{l-1}}{c_{\text{out}} c_{\text{in}} k_h k_w}
\end{equation}
where $c_{\text{out}}, c_{\text{in}}$ are the channels, and $k_h, k_w$ are kernel dimensions. During the forward pass of the test batch, we register forward hooks to save input activations. During the backward pass of the entropy loss, backward hooks capture the loss gradients. The sensitivities $\bar{F}_w$ are then computed dynamically in a single pass.

\subsection{The Stabilizer Drowning Effect and Global Normalization}
A critical issue identified in our analysis of standard on-the-fly KFAC trace estimation is the \textbf{stabilizer drowning effect}. Raw sensitivity values computed from $\bar{F}_w$ are divided by the number of parameter connections, making their absolute magnitudes extremely small (often on the order of $10^{-8}$ or $10^{-10}$). When these sensitivities are used in preconditioned gradient updates:
\begin{equation}
    \delta_j^{(t+1)} = \delta_j^{(t)} - \eta \cdot \frac{1}{\bar{F}_j + \epsilon_{\text{stab}}} \cdot \frac{\partial L}{\partial \delta_j}
\end{equation}
with standard stabilizers like $\epsilon_{\text{stab}} = 10^{-2}$ or even $10^{-4}$, the sensitivity $\bar{F}_j$ is completely dominated by $\epsilon_{\text{stab}}$:
\begin{equation}
    \bar{F}_j + \epsilon_{\text{stab}} \approx \epsilon_{\text{stab}}
\end{equation}
This effectively neutralizes the layer-wise preconditioning, making the update scales identical across all layers. To resolve this, we introduce a \textbf{Global Normalization} strategy. We track all estimated layer-wise sensitivities across the network and divide each of them by the maximum sensitivity value in the current pass:
\begin{equation}
    \hat{F}_j = \frac{\bar{F}_j}{\max_{k} \bar{F}_k}
\end{equation}
This scales the sensitivity values to the interval $[0, 1]$, making them directly comparable to $\epsilon_{\text{stab}}$ and fully restoring the preconditioning characteristics of on-the-fly Kronecker-Fisher.

\subsection{Momentum-Augmented (EMA) Sensitivity Tracking}
While on-the-fly Kronecker trace sensitivity estimation operates zero-shot and data-free, single-batch estimates can be inherently noisy, especially under non-stationary streams with small batch sizes or transient outliers. Sudden fluctuations in sensitivity values can lead to erratic preconditioning, causing unstable adaptation or representation drift in highly sensitive layers. 

To address this, we propose \textbf{Momentum-Augmented Kronecker-Fisher preconditioning}, which smooths on-the-fly sensitivities across the streaming horizon. We maintain a running Exponential Moving Average (EMA) of the globally-normalized sensitivities:
\begin{equation}
    \tilde{F}_j^{(t)} = \alpha_{\text{ema}} \tilde{F}_j^{(t-1)} + (1 - \alpha_{\text{ema}}) \hat{F}_j^{(t)}
\end{equation}
where $\alpha_{\text{ema}} \in [0, 1)$ is the smoothing momentum parameter, and $\hat{F}_j^{(t)}$ is the globally-normalized sensitivity computed on the current test batch. For the initial batch ($t=0$), we initialize the running state as $\tilde{F}_j^{(0)} = \hat{F}_j^{(0)}$. The preconditioned update is then computed using the smoothed sensitivity $\tilde{F}_j^{(t)}$:
\begin{equation}
    \delta_j^{(t+1)} = \delta_j^{(t)} - \eta \cdot \frac{1}{\tilde{F}_j^{(t)} + \epsilon_{\text{stab}}} \cdot \frac{\partial L}{\partial \delta_j}
\end{equation}
This momentum-augmented formulation retains historical structural awareness of parameter sensitivities from prior batches along the stream, stabilizing adaptation and preventing representational drift under noise or domain transitions.

\subsection{Variance-Preserving Soft BN Buffer Fusion}
Merging model weights alters the intermediate activation distributions, causing a mismatch with the running statistics of the Batch Normalization (BN) layers. Simple linear interpolation of mean and variance fails to model the mixture distribution. Following the law of total variance, we implement \textit{Variance-Preserving Soft BN Fusion} to reconstruct exact mixture moments:
\begin{align}
    \mu_{\text{fused}} &= w_1 \mu_1 + w_2 \mu_2 \\
    \sigma_{\text{fused}}^2 &= w_1 \sigma_1^2 + w_2 \sigma_2^2 \nonumber \\
    &\quad + w_1(\mu_1 - \mu_{\text{fused}})^2 + w_2(\mu_2 - \mu_{\text{fused}})^2
\end{align}
where $w_1 = \bar{\lambda}$ and $w_2 = 1 - \bar{\lambda}$ are the average mixture weights across all layers. This preserves activation magnitudes, preventing vanishing or exploding activations in deep layers. However, in simple grayscale domains with highly isolated and clean activations, we find that simple linear moment averaging may prevent overestimating variance from mean differences, a trade-off we explore in our ablations.

\subsection{KL-Regularized Test-Time Adaptation}
When a test batch $X^{(t)}$ arrives, we project features into the Static Unified Space of each expert and compute the similarity scores $S_k(X^{(t)})$ to class prototypes. Under Self-Calibrated Temperature Scaling (SCTS) \citep{yang2025testtime}, the temperature $\tau$ is scaled by the prototype distance gap $\Delta$:
\begin{equation}
    \tau = \frac{\Delta}{s} + \epsilon_{\text{stab}}
\end{equation}
The routing prior probability is $p = \text{softmax}(S/\tau)_1$. We initialize the global logit as $w_{\text{global}}^{(0)} = \log(p / (1-p))$ and offsets $\delta_j^{(0)} = 0.0$. We then perform test-time adaptation to minimize prediction entropy, regularized by the KL divergence to the routing prior and the offset coherence penalty:
\begin{equation}
    L = L_{\text{entropy}} + \beta D_{\text{KL}}(w \| \lambda) + L_{\text{coherence}}
\end{equation}
Updates are preconditioned by the globally normalized Kronecker-Fisher sensitivities $\hat{F}_j$:
\begin{align}
    w_{\text{global}}^{(t+1)} &= w_{\text{global}}^{(t)} - \eta \cdot \frac{\partial L}{\partial w_{\text{global}}} \\
    \delta_j^{(t+1)} &= \delta_j^{(t)} - \eta \cdot \frac{1}{\hat{F}_j + \epsilon_{\text{stab}}} \cdot \frac{\partial L}{\partial \delta_j}
\end{align}

\subsection{Stability Analysis of Co-acting Layer-Wise Offsets}
To theoretically justify how the co-acting layer-wise offset parameterization prevents representation collapse during unconstrained test-time adaptation, we analyze the discrete-time dynamical system governing the offset updates. We prove that under a suitable learning rate condition, the layer-wise offsets remain stable and bounded, preventing individual layers from drifting arbitrarily far from the global consensus logit.

\begin{theorem}[Stability of Co-acting Offsets]
\label{thm:stability}
Assume that the gradient of the unsupervised task-adaptation loss with respect to the layer-wise offsets is bounded, i.e., $\|\nabla_{\delta_j} L_{\text{entropy}}\|_2 \le G_{\max}$ for all layers $j$ and iterations $t$. Let $\hat{F}_j \in [0, 1]$ be the globally-normalized sensitivity, and let $\epsilon_{\text{stab}} > 0$ be the preconditioning stabilizer. If the learning rate $\eta$ satisfies the stability condition:
\begin{equation}
    \eta \le \frac{\hat{F}_j + \epsilon_{\text{stab}}}{2\gamma}
\end{equation}
then the layer-specific offset $\delta_j^{(t)}$ is stable and remains bounded for all $t \ge 0$. Specifically, starting from the zero initialization $\delta_j^{(0)} = 0.0$, the asymptotic bound is:
\begin{equation}
    \limsup_{t \to \infty} \|\delta_j^{(t)}\|_2 \le \frac{G_{\max}}{2\gamma}
\end{equation}
\end{theorem}

\begin{proof}
Let $\rho_j = 1 - \frac{2\gamma\eta}{\hat{F}_j + \epsilon_{\text{stab}}}$. Applying the update rule for $\delta_j^{(t+1)}$, we have:
\begin{equation*}
    \delta_j^{(t+1)} = \rho_j \delta_j^{(t)} - \frac{\eta}{\hat{F}_j + \epsilon_{\text{stab}}} \nabla_{\delta_j} L_{\text{entropy}}^{(t)}
\end{equation*}
Taking the L2 norm on both sides and applying the triangle inequality yields:
\begin{align*}
    \|\delta_j^{(t+1)}\|_2 &\le |\rho_j| \|\delta_j^{(t)}\|_2 \\
    &\quad + \frac{\eta}{\hat{F}_j + \epsilon_{\text{stab}}} \|\nabla_{\delta_j} L_{\text{entropy}}^{(t)}\|_2
\end{align*}
Under the learning rate stability condition $\eta \le \frac{\hat{F}_j + \epsilon_{\text{stab}}}{2\gamma}$, we have:
\begin{equation*}
    0 < \frac{2\gamma\eta}{\hat{F}_j + \epsilon_{\text{stab}}} \le 1 \implies 0 \le \rho_j < 1
\end{equation*}
Thus, $|\rho_j| = \rho_j \in [0, 1)$. Using the bounded gradient assumption $\|\nabla_{\delta_j} L_{\text{entropy}}^{(t)}\|_2 \le G_{\max}$, we obtain the recurrence relation:
\begin{equation*}
    \|\delta_j^{(t+1)}\|_2 \le \rho_j \|\delta_j^{(t)}\|_2 + \frac{\eta G_{\max}}{\hat{F}_j + \epsilon_{\text{stab}}}
\end{equation*}
By induction, starting from $\delta_j^{(0)} = 0.0$, we have:
\begin{align*}
    \|\delta_j^{(t)}\|_2 &\le \left( \sum_{k=0}^{t-1} \rho_j^k \right) \frac{\eta G_{\max}}{\hat{F}_j + \epsilon_{\text{stab}}} \\
    &= \frac{1 - \rho_j^t}{1 - \rho_j} \frac{\eta G_{\max}}{\hat{F}_j + \epsilon_{\text{stab}}}
\end{align*}
Taking the limit as $t \to \infty$, since $0 \le \rho_j < 1$:
\begin{equation*}
    \limsup_{t \to \infty} \|\delta_j^{(t)}\|_2 \le \frac{1}{1 - \rho_j} \frac{\eta G_{\max}}{\hat{F}_j + \epsilon_{\text{stab}}}
\end{equation*}
Substituting $1 - \rho_j = \frac{2\gamma\eta}{\hat{F}_j + \epsilon_{\text{stab}}}$ back into the inequality:
\begin{align*}
    \limsup_{t \to \infty} \|\delta_j^{(t)}\|_2 &\le \left( \frac{\hat{F}_j + \epsilon_{\text{stab}}}{2\gamma\eta} \right) \frac{\eta G_{\max}}{\hat{F}_j + \epsilon_{\text{stab}}} \\
    &= \frac{G_{\max}}{2\gamma}
\end{align*}
This completes the proof.
\end{proof}

Our complete execution logic is detailed in Algorithm \ref{alg:cl_kt_fisher}.

\begin{algorithm}[tb]
\caption{CL-KT-Fisher Adaptation Loop}
\label{alg:cl_kt_fisher}
\begin{algorithmic}[1]
   \STATE {\bfseries Input:} Test stream batch $X^{(t)}$, expert weights $\theta_1, \theta_2$, class prototypes, learning rate $\eta$, regularizations $\beta, \gamma$.
   \STATE Extract features and compute SCTS routing prior $p_{\text{prior}}$.
   \STATE Initialize $w_{\text{global}} \leftarrow \log(p_{\text{prior}} / (1 - p_{\text{prior}}))$, and $\delta_j \leftarrow 0.0$ for all layers $j$.
   \STATE Register Kronecker Trace forward and backward hooks.
   \STATE Compute single-step backward pass of prediction entropy on $X^{(t)}$.
   \STATE Extract and globally normalize sensitivities: $\hat{F}_j \leftarrow \frac{\bar{F}_j}{\max_k \bar{F}_k}$.
   \STATE Remove hooks and clear temporary gradients.
   \STATE Compute combined loss: $L \leftarrow L_{\text{entropy}} + \beta D_{\text{KL}}(w \| \lambda) + L_{\text{coherence}}$.
   \STATE Compute gradients of $L$ with respect to $w_{\text{global}}$ and all $\delta_j$.
   \STATE Update global logit: $w_{\text{global}} \leftarrow w_{\text{global}} - \eta \frac{\partial L}{\partial w_{\text{global}}}$.
   \STATE Update offsets: $\delta_j \leftarrow \delta_j - \eta \frac{1}{\hat{F}_j + 10^{-2}} \frac{\partial L}{\partial \delta_j}$.
   \STATE Merge parameters differentiably to evaluate and predict.
\end{algorithmic}
\end{algorithm}

\section{Experiments}
\label{sec:experiments}

\subsection{Experimental Setup}
We train two ResNet-18 expert models on MNIST and KMNIST from scratch using AdamW (learning rate $10^{-3}$, weight decay $10^{-4}$) for 3 epochs. We designate FashionMNIST as a completely novel domain. This simulates open-world streams where novel task distributions emerge dynamically. Standalone expert test accuracies are: MNIST Expert on MNIST: \textbf{98.83\%}, KMNIST Expert on KMNIST: \textbf{94.32\%}, and FashionMNIST Expert on FashionMNIST: \textbf{90.28\%}.

We construct four evaluation stream profiles consisting of batches of size 64:
\begin{itemize}
    \item \textbf{Profile 1 (Closed Sequential):} 15 batches of MNIST, followed by 15 batches of KMNIST.
    \item \textbf{Profile 2 (Closed Alternating):} 30 alternating batches of MNIST and KMNIST.
    \item \textbf{Profile 3 (Open-World Short):} 10 batches MNIST, 10 batches KMNIST, and 10 batches FashionMNIST (novel).
    \item \textbf{Profile 4 (Open-World Standard):} 30 batches MNIST, 30 batches KMNIST, and 30 batches FashionMNIST (novel).
\end{itemize}

\subsection{Baselines}
We compare our proposed \textbf{CL-KT-Fisher} against five strong baselines:
\begin{itemize}
    \item \textbf{Static Uniform:} Merges expert parameters with a static average weight of 0.5 without any test-time adaptation.
    \item \textbf{AdaMerging (Standard):} Dynamically adapts a single global coefficient via unconstrained test-time entropy minimization.
    \item \textbf{KT-Fisher (Paper 9):} Uses Kronecker trace-based preconditioning but restricts updates to simple, layer-naive scalar scaling.
    \item \textbf{CL W-Fisher (Paper 8):} Uses the global logit + local offset co-acting parameterization preconditioned by \textit{offline precomputed} Joint Fisher Information on calibration datasets.
    \item \textbf{DF-Bayes-TTMM (Paper 10):} Implements Bayesian posterior routing and Soft BN fusion but operates without prior-guided parameter initialization.
\end{itemize}

\section{Results}
\label{sec:results}

\subsection{Performance under Closed-Loop Streams}
As shown in the overall accuracy comparison in \cref{fig:accuracy_comparison} and detailed segment-level results in \cref{tab:results_profile_1_2}, Static Uniform, AdaMerging, and KT-Fisher struggle on Profile 1 and Profile 2 (scoring around 45.47\% overall). Standard unconstrained entropy minimization methods quickly fall into the ``feedback loop trap,'' collapsing representations. DF-Bayes-TTMM also performs poorly (41.41\%) without proper prior-guided weight initialization. 

\begin{figure}[htbp]
\centering
\includegraphics[width=\columnwidth]{figures/accuracy_comparison}
\caption{Overall stream accuracy comparison of various model merging methods on Profile 1 (Closed Sequential) and Profile 4 (Open-World Standard).}
\label{fig:accuracy_comparison}
\end{figure}

In contrast, our proposed \textbf{CL-KT-Fisher (Ours)} achieves a remarkable accuracy of \textbf{86.88\%} on both profiles. It closely matches the performance of the data-dependent CL W-Fisher baseline (90.42\%) while being \textbf{entirely data-free}, requiring no access to offline calibration sets. This demonstrates the immense value of combining co-acting parameterization with online Kronecker-Fisher preconditioning.

\begin{table*}[t]
\centering
\caption{Segment-level and overall stream accuracies (\%) on Profile 1 (Closed Sequential) and Profile 2 (Closed Alternating).}
\label{tab:results_profile_1_2}
\vskip 0.15in
\begin{small}
\setlength{\tabcolsep}{3.5pt}
\begin{tabular}{lccccccc}
\toprule
& & \multicolumn{3}{c}{\textbf{Profile 1 (Closed Sequential)}} & \multicolumn{3}{c}{\textbf{Profile 2 (Closed Alternating)}} \\
\cmidrule(lr){3-5} \cmidrule(lr){6-8}
Method & DF & MNIST & KMNIST & \textbf{Overall} & MNIST & KMNIST & \textbf{Overall} \\
\midrule
Static Uniform & \checkmark & 67.08 & 23.54 & 45.31 & 67.08 & 23.54 & 45.31 \\
AdaMerging (Standard) & \checkmark & 67.40 & 23.54 & 45.47 & 67.40 & 23.54 & 45.47 \\
KT-Fisher (Paper 9) & \checkmark & 67.40 & 23.54 & 45.47 & 67.40 & 23.54 & 45.47 \\
DF-Bayes-TTMM (Paper 10) & \checkmark & 61.98 & 20.83 & 41.41 & 61.98 & 20.83 & 41.41 \\
\midrule
CL W-Fisher (Paper 8) & \textsf{X} & \textbf{98.75} & \textbf{82.08} & \textbf{90.42} & \textbf{98.75} & \textbf{82.08} & \textbf{90.42} \\
\textbf{CL-KT-Fisher (Ours)} & \checkmark & 98.33 & 75.42 & \underline{86.88} & 98.33 & 75.42 & \underline{86.88} \\
\ \ \textbf{Ours (No VP-BN)} & \checkmark & \textbf{98.75} & \textbf{82.08} & \textbf{90.42} & \textbf{98.75} & \textbf{82.08} & \textbf{90.42} \\
\ \ \textbf{Ours (EMA)} & \checkmark & 98.33 & 75.42 & \underline{86.88} & 98.33 & 75.42 & \underline{86.88} \\
\ \ \textbf{Ours (EMA, No VP-BN)} & \checkmark & \textbf{98.75} & \textbf{82.08} & \textbf{90.42} & \textbf{98.75} & \textbf{82.08} & \textbf{90.42} \\
\bottomrule
\end{tabular}
\end{small}
\end{table*}

\subsection{Performance under Open-World Streams}
In open-world streams containing the completely unseen FashionMNIST domain (Profiles 3 and 4), our proposed \textbf{CL-KT-Fisher (Ours)} demonstrates superior robustness. As shown in \cref{tab:results_profile_3_4}, it achieves \textbf{61.88\%} on Profile 3 and \textbf{62.27\%} on Profile 4, significantly outperforming layer-naive methods (which score 32.92\% and 34.55\% respectively). 

Notably, on the completely novel FashionMNIST task under Profile 4, CL-KT-Fisher achieves a standalone accuracy of \textbf{12.92\%}, outperforming even the data-dependent CL W-Fisher (12.71\%) and Static Uniform (12.03\%). This indicates that dynamic, online sensitivity estimation allows the model to protect and utilize robust expert pathways even on tasks it was never trained on.

\begin{table*}[t]
\centering
\caption{Segment-level and overall stream accuracies (\%) on Profile 3 (Open-World Short) and Profile 4 (Open-World Standard).}
\label{tab:results_profile_3_4}
\vskip 0.15in
\begin{footnotesize}
\setlength{\tabcolsep}{3.0pt}
\begin{tabular}{lccccccccc}
\toprule
& & \multicolumn{4}{c}{\textbf{Profile 3 (Open-World Short)}} & \multicolumn{4}{c}{\textbf{Profile 4 (Open-World Standard)}} \\
\cmidrule(lr){3-6} \cmidrule(lr){7-10}
Method & DF & MNIST & KMNIST & Fashion & \textbf{Overall} & MNIST & KMNIST & Fashion & \textbf{Overall} \\
\midrule
Static Uniform & \checkmark & 67.03 & 21.56 & 9.84 & 32.81 & 67.86 & 23.59 & 12.03 & 34.50 \\
AdaMerging & \checkmark & 67.34 & 21.56 & 9.84 & 32.92 & 68.02 & 23.59 & 12.03 & 34.55 \\
KT-Fisher (Paper 9) & \checkmark & 67.34 & 21.56 & 9.84 & 32.92 & 68.02 & 23.59 & 12.03 & 34.55 \\
DF-Bayes-TTMM & \checkmark & 61.56 & 18.59 & 10.00 & 30.05 & 61.09 & 20.47 & 11.98 & 31.18 \\
\midrule
CL W-Fisher (Paper 8) & \textsf{X} & \textbf{98.91} & \textbf{81.25} & 10.94 & \textbf{63.70} & 98.65 & \textbf{81.93} & 12.71 & \textbf{64.43} \\
\textbf{CL-KT-Fisher (Ours)} & \checkmark & 98.44 & 75.62 & \textbf{11.56} & \underline{61.88} & 98.28 & 75.62 & \textbf{12.92} & \underline{62.27} \\
\ \ \textbf{Ours (No VP-BN)} & \checkmark & \textbf{98.91} & \textbf{81.25} & 10.94 & \textbf{63.70} & \textbf{98.65} & \textbf{81.93} & 12.71 & \textbf{64.43} \\
\ \ \textbf{Ours (EMA)} & \checkmark & 98.44 & 75.62 & \textbf{11.56} & \underline{61.88} & 98.28 & 75.62 & \textbf{12.92} & \underline{62.27} \\
\ \ \textbf{Ours (EMA, No VP-BN)} & \checkmark & \textbf{98.91} & \textbf{81.25} & 10.94 & \textbf{63.70} & \textbf{98.65} & \textbf{81.93} & 12.71 & \textbf{64.43} \\
\bottomrule
\end{tabular}
\end{footnotesize}
\end{table*}

\subsection{Ablation and Hyperparameter Analysis}
We conduct a detailed ablation study on Profile 1 (Closed Sequential) and Profile 4 (Open-World Standard) to analyze the impact of our core design components and understand the merging dynamics under streaming shifts (visualized in \cref{fig:coefficient_trajectory}):

\begin{figure}[htbp]
\centering
\includegraphics[width=\columnwidth]{figures/coefficient_trajectory}
\caption{Trajectory of the merging coefficient $\lambda$ over 30 test-time batches of Profile 1. CL-KT-Fisher (Ours) successfully tracks the ideal prior-guided trajectory dynamically without any offline pretraining, whereas Standard AdaMerging is unable to adapt due to representational collapse.}
\label{fig:coefficient_trajectory}
\end{figure}

\begin{enumerate}
    \item \textbf{Role of Variance-Preserving BN Fusion:} We compare our exact moment-based Soft BN Fusion (Ours) against simple linear interpolation of expert BN running statistics (\textsf{No VP-BN}). On these grayscale vision experts, simple linear interpolation achieves slightly higher accuracies (90.42\% on Profile 1 and 64.43\% on Profile 4), matching the data-dependent CL W-Fisher oracle. This suggests that for clean, isolated activations, simple linear moment averaging preserves activation scales perfectly, avoiding overestimating variance from mean differences.
    \item \textbf{Sensitivity of Coherence Regularization ($\gamma$):} We sweep the consensus coherence penalty $\gamma \in \{0.0, 0.002, 0.02, 0.1, 0.5\}$ to evaluate the stability of our co-acting representation. Across all sweeps, our method maintains a constant accuracy of 86.88\% on Profile 1 and 62.27\% on Profile 4. This extreme robustness to hyperparameter selection is a major advantage for unmonitored test-time deployments where tuning is impossible.
    \item \textbf{Sensitivity of KL Regularization ($\beta$):} We sweep the KL regularization parameter $\beta \in \{0.1, 1.0, 5.0\}$ to test the influence of the prior-guided regularization strength. Across all sweeps, our method achieves exactly identical overall accuracies. This highlights that our SCTS-based prior-guided routing and initialization is so highly optimal that the exact weight of the KL regularization penalty has negligible influence on the local optimization landscape, ensuring superb scientific stability.
    \item \textbf{Optimization Learning Rate ($\eta$):} We evaluate the impact of different learning rates $\eta \in \{0.005, 0.01, 0.05\}$ on adaptation. We find that a larger learning rate $\eta = 0.05$ slightly improves performance (\textbf{87.03\%} on Profile 1 and \textbf{62.48\%} on Profile 4) by facilitating faster adaptation within a single step. Conversely, a smaller learning rate of $0.005$ slightly dampens adaptation, leading to $86.67\%$ and $62.22\%$ respectively. This confirms that CL-KT-Fisher is highly robust to learning rate selection within an order of magnitude.
    \item \textbf{Stability of On-the-Fly Sensitivity Profiles (EMA Tracking):} We evaluate the effect of the Momentum-Augmented Exponential Moving Average (EMA) sensitivity tracking (\textsf{EMA} variants) with $\alpha_{\text{ema}} = 0.8$. Strikingly, we find that both the standard and EMA variants achieve identical overall accuracies across all stream profiles (e.g., 86.88\% on Profile 1 and 62.27\% on Profile 4). This indicates that the relative layer-wise structural sensitivity profiles are highly stable characteristics of the ResNet-18 architecture under different domain shifts. Consequently, the normalized preconditioning matrix remains consistent from batch to batch, rendering our framework remarkably robust to transient noise without requiring historical momentum smoothing.
\end{enumerate}

\begin{table}[t]
\centering
\caption{Ablation study and hyperparameter sweeps on Profile 1 and Profile 4 (Accuracies \%). Default configuration (Ours) uses $\eta=0.01$, $\beta=1.0$, and $\gamma=0.02$.}
\label{tab:ablations}
\vskip 0.15in
\begin{small}
\setlength{\tabcolsep}{5pt}
\begin{tabular}{lcc}
\toprule
Variant & Profile 1 & Profile 4 \\
\midrule
\textbf{CL-KT-Fisher (Ours)} & 86.88 & 62.27 \\
\ \ \textsf{No VP-BN} & \textbf{90.42} & \textbf{64.43} \\
\ \ \textsf{EMA} & 86.88 & 62.27 \\
\ \ \textsf{EMA, No VP-BN} & \textbf{90.42} & \textbf{64.43} \\
\midrule
\textit{Learning Rate Sweep} ($\eta$) & & \\
\ \ $\eta = 0.05$ & \underline{87.03} & \underline{62.48} \\
\ \ $\eta = 0.005$ & 86.67 & 62.22 \\
\midrule
\textit{KL Regularization Sweep} ($\beta$) & & \\
\ \ $\beta = 0.1$ & 86.88 & 62.27 \\
\ \ $\beta = 5.0$ & 86.88 & 62.27 \\
\midrule
\textit{Coherence Penalty Sweep} ($\gamma$) & & \\
\ \ $\gamma = 0.0$ & 86.88 & 62.27 \\
\ \ $\gamma = 0.002$ & 86.88 & 62.27 \\
\ \ $\gamma = 0.1$ & 86.88 & 62.27 \\
\ \ $\gamma = 0.5$ & 86.88 & 62.27 \\
\bottomrule
\end{tabular}
\end{small}
\end{table}

\section{Discussion and Future Work}
\label{sec:discussion}
While CL-KT-Fisher achieves excellent results on non-stationary streams, several theoretical and empirical properties merit further discussion. First, the success of the data-free Kronecker trace-based preconditioning depends on the quality of the on-the-fly estimated gradients of prediction entropy. Because we use a single forward/backward pass per batch, small batch sizes can introduce noise in KFAC trace estimates, though we find a batch size of 64 provides highly stable estimates. Second, the slight performance gap of Variance-Preserving Soft BN Fusion compared to simple linear interpolation highlights a fundamental assumption mismatch in streaming scenarios. The variance-preserving formula based on the law of total variance incorporates a variance-inflation term proportional to the squared difference in means, $w_1(1-w_1)(\mu_1-\mu_2)^2$, which assumes that the incoming test batch is a mixture of samples from both domains. However, in our non-stationary sequential stream, each test batch is homogeneous (containing samples purely from a single domain). When evaluating such a homogeneous batch, this inflation term overestimates the actual variance of the batch, causing over-normalization of activations in deep layers. This over-normalization dampens the representation signal, leading to a slight drop in accuracy compared to simple linear interpolation ($w_1\sigma_1^2 + w_2\sigma_2^2$), which does not suffer from this overestimation effect. This indicates that variance-preserving fusion is best suited for heterogeneous batches, whereas standard linear interpolation is preferred for sequential, homogeneous batch streams.

Future work includes scaling CL-KT-Fisher to large-scale foundation models (such as LLMs and vision-language models), where layer sensitivities are even more diverse and offline data dependencies are more severe. Additionally, investigating exponential moving averages of Kronecker trace sensitivities across a sliding test window could provide smoother adaptation trajectories.

\section{Conclusion}
\label{sec:conclusion}
In this work, we introduced Co-acting Layer-Wise Kronecker-Fisher Adaptation (CL-KT-Fisher) to address the core challenges of representation collapse, layer-wise parameter sensitivity, and data dependencies in online Test-Time Model Merging. By utilizing an expressive co-acting parameterization and estimating Kronecker trace sensitivities on-the-fly directly from test streams, CL-KT-Fisher achieves true data-free test-time model merging. Our extensive evaluation shows that our method achieves outstanding accuracies across multiple non-stationary streams, representing a significant advancement towards robust, zero-shot open-world model deployments.

\nocite{marczak2025task, wu2025unlocking, chen2025bring, cheng2025whoever, wei2025modeling, li2025model, tang2025merging, corbeil2025modular, yuan2025superficial, yang2025data, gargiulo2024task, panariello2025accurate, du2025adamms, kuo2025exact, sun2025merging, akiba2024evolutionary, yang2024model, wei2025disrupting, lu2024finetuning, niu2023towards, karmanov2024efficient, yuan2023robust, chen2022contrastive, lee2024entropy, kim2025battling, guo2025smoothing, zhou2025bayesian, guo2025everything, cui2025online, zhang2025multisource, song2023ecotta, boudiaf2022parameterfree, dbler2022robust, liang2023comprehensive, gong2022note, han2025ranked, dong2025towards, zanella2025realistic, liu2025fedswa, yu2025sample}

\bibliography{example_paper}
\bibliographystyle{icml2026}

\end{document}
